\section{Global classification}\label{sec:global}

We now prove the main equivalence.

\begin{proof}[Proof of \cref{thm:main}: necessity]
Assume \(N\preceqplus N'\).  The noncircular cut-transfer theorem
\cref{thm:cut-transfer} gives equality of the labelled bridge splits and hence
the same labelled reduced tree of blobs.  The complete bridge fibre and its
physical product chart extract corresponding projective local factors with
exactly three incidence gauges and no holonomy.  Physical localization
\cref{prop:localization} fixes one target incoming/completion type for each
source-open focal germ and excludes remote compensation.  The complete local
classification \cref{thm:local} then makes every corresponding factor either
labelled-isomorphic or ordinary-triangle related.  The one-/two-port
transports reconstruct every segment word and make all triangle choices
coherent.  Thus \(N\trianglerel N'\).
\end{proof}

It remains to realize the structural relation physically and at full local
dimension.

\begin{lemma}[Simultaneous physical bridge gluing]\label{lem:gluing}
Suppose corresponding local factors have one common projective K3P germ,
through isomorphism or \cref{lem:H14-context}.  After shrinking finitely many
local germs, their contractions along the common bridge tree form a common
full-dimensional regular global germ, both on \(\Dplus\) and on \(\Dct\).
\end{lemma}

\begin{proof}
Relative to fixed analytic bridge slices, let
\(A_h=a^h_{u,e}a^h_{v,e}\) be the endpoint incidence product on either network
and any bridge, with \(h=C,G,T\).  On compactly contained local germs choose
uniform bounds \(0<L\le A_h\le U\).  Put

\[
 \varepsilon=\frac{L^2}{4U},\qquad
 z_C=z_G=z_T=\varepsilon,\qquad
 x_h=\frac{\varepsilon}{A_h}.
\]

Then

\[
 \frac{L^2}{4U^2}\le x_h\le\frac{L}{4U}\le\frac14.
\]

Every principal-domain composition margin is at least
\(1-L/(2U)\ge1/2\), and every continuous-time margin is at least

\[
 \frac{L^2}{4U^2}-\frac{L^2}{16U^2}
 =\frac{3L^2}{16U^2}>0.
\]

Thus the same effective isotropic bridge spectrum works simultaneously for
both networks and all bridges.  Incidence factors cancel in
\eqref{eq:bridge-contraction}.  The physical product extraction in
\cref{lem:physical-product} is a local inverse, so the contraction has the
expected rank and the common image germ is full-dimensional.
\end{proof}

\begin{proof}[Proof of \cref{thm:main}: sufficiency]
Assume \(N\trianglerel N'\).  Transport physical parameters across every
labelled isomorphism.  For each redirected triangle use the same smooth
relative \(H_{14}\) germ and its physical analytic sections; choose the
generic rank once in the joint contextual contraction, as in
\cref{lem:H14-context}.  Apply \cref{lem:gluing} to all bridges
simultaneously.  The resulting common germ is regular and full-dimensional
in both complete network images, so \(N\bowtieplus N'\).  A common germ with
physical sections gives the directed relation in both directions by
definition.  This proves \eqref{eq:main-classification}.  Since
\(\trianglerel\) is symmetric, proper one-sided containment is impossible.
\end{proof}

The triangle part of this proof is full-dimensional \emph{relative to each
network image}.  No step replaces \(H_{14}\) by an ambient-open
\(15\)-dimensional triangle model.
